\section{Metodologia}
\label{sec:metodologia}

\subsection{Pipeline experimental em seis etapas}
\label{sec:pipeline}

O pipeline metodológico aprovado pelo orientador estrutura-se em seis etapas sequenciais:

\begin{enumerate}
\item \textbf{Etapa 1 --- Classificador de tom de pele MST}: usar SkinToneNet pré-treinado \autocite{pereira2026skintonenet} como insumo, mediante validação independente em subset FairFace anotado manualmente.

\item \textbf{Etapa 2 --- Auditoria FairFace}: aplicar SkinToneNet sobre FairFace validation set e construir a matriz pública MST $\times$ race classes (Contribuição C2).

\item \textbf{Etapa 3 --- Race classifier com conditioning}: treinar ConvNeXt-T fine-tuned em FairFace train, com camadas FiLM por estágio recebendo vetor MST como contexto.

\item \textbf{Etapa 4 --- Comparação contra baselines}: avaliar o pipeline contra seis baselines de mitigação, reportando triangulação completa de métricas (Contribuição C4).

\item \textbf{Etapa 5 --- Fair transferência}: aplicar o backbone fair em face recognition downstream (RFW ou BFW), com controle de pixel information.

\item \textbf{Etapa 6 --- Síntese decompositiva}: combinar resultados de C2 e C5 para quantificar componente fenotípico versus algorítmico do erro Latinx (Contribuição C6).
\end{enumerate}

\subsection{Datasets}
\label{sec:datasets}

\begin{description}[font=\bfseries\sffamily]
\item[Treino race]: FairFace train split (86.744 imagens, sete classes raciais) \autocite{karkkainen2021fairface}.
\item[Validação MST]: STW \autocite{pereira2026skintonenet} e Casual Conversations v2 \autocite{porgali2023ccv2}.
\item[Teste race]: FairFace test split (10.954 imagens) \autocite{karkkainen2021fairface}.
\item[Teste fair transfer]: RFW \autocite{wang2019rfw} e/ou BFW \autocite{robinson2020bfw}.
\end{description}

\subsection{Baselines}
\label{sec:baselines}

Comparação contra seis baselines de mitigação algorítmica:
\begin{enumerate}
\item ResNet-34 (canônico FairFace) \autocite{karkkainen2021fairface,aldahoul2026};
\item ConvNeXt-T puro (controle arquitetural moderno);
\item FSCL+ \autocite{park2022fscl};
\item Group DRO \autocite{sagawa2020groupdro};
\item FineFACE \autocite{manzoor2024fineface};
\item Adversarial debiasing \autocite{zhang2018adversarial}.
\end{enumerate}

\subsection{Mecanismo de conditioning FiLM}
\label{sec:film}

Para cada estágio do ConvNeXt-T, insere-se uma camada FiLM \autocite{perez2018film} que recebe o vetor MST $z \in \mathbb{R}^{10}$ (saída softmax do SkinToneNet pré-treinado) e produz parâmetros de modulação $(\gamma_i, \beta_i)$ por canal, aplicados como transformação afim:

\begin{equation}
\mathrm{FiLM}(F_i \mid \gamma_i, \beta_i) = \gamma_i \odot F_i + \beta_i
\end{equation}

onde $F_i$ é a feature map intermediária do estágio $i$ e $\odot$ denota multiplicação elemento-a-elemento. As funções $f_\gamma$ e $f_\beta$ são duas pequenas MLPs treináveis com aproximadamente 380.000 parâmetros adicionais ($\approx 1\%$ de overhead sobre o backbone de 28 milhões de parâmetros).

\subsection{Ablation arquitetural --- FiLM versus CLIP-conditioning}
\label{sec:ablation}

Endereçando a recomendação registrada na reunião de orientação de 2026-06-08, será conduzida ablation arquitetural comparando FiLM \autocite{perez2018film} contra abordagem CLIP-based \autocite{luo2024fairclip,radford2021clip}, esta última empregando embeddings de CLIP pré-treinado como contexto alternativo via cross-attention.

\subsection{Triangulação de métricas}
\label{sec:metricas}

Fundamentada no Teorema da Impossibilidade \autocite{kleinberg2017}, a triangulação envolve dois cenários:

\paragraph{Cenário A --- Protocolo principal (race apenas).}
\begin{itemize}
\item F1 macro \autocite{vanrijsbergen1979};
\item Disparity Ratio: $\mathrm{DR} = \min_c \mathrm{F1}_c / \max_c \mathrm{F1}_c$ \autocite{hardt2016eqopp};
\item Worst-class F1 \autocite{sagawa2020groupdro}.
\end{itemize}

\paragraph{Cenário B --- Ablation intersectional (race $\times$ gender).}
\begin{itemize}
\item Equal Opportunity por classe: para cada raça $c$, gap entre $\mathrm{TPR}_c^{M}$ e $\mathrm{TPR}_c^{F}$ \autocite{hardt2016eqopp};
\item Equalized Odds por classe: gap simultâneo em TPR e FPR \autocite{hardt2016eqopp}.
\end{itemize}

\subsection{Rigor estatístico}
\label{sec:rigor}

\begin{itemize}
\item Três seeds independentes (42, 1, 2) com média $\pm$ desvio padrão;
\item Comparação pareada contra baselines;
\item Intervalos de confiança via bootstrap.
\end{itemize}
